% Generated draft: requires mathematical and source review.
\documentclass[11pt]{article}
\usepackage[margin=25mm]{geometry}
\usepackage{fontspec}
\setmainfont{Latin Modern Roman}
\usepackage{amsmath,amssymb,mathtools,unicode-math}
\setmathfont{Latin Modern Math}
\usepackage{enumitem,needspace,xurl,hyperref}
\setlength{\parindent}{0pt}
\setlength{\parskip}{5pt}
\setlength{\emergencystretch}{3em}
\newcommand{\N}{\mathbb{N}}
\newcommand{\Z}{\mathbb{Z}}
\newcommand{\Q}{\mathbb{Q}}
\newcommand{\R}{\mathbb{R}}
\newcommand{\C}{\mathbb{C}}
\newcommand{\F}{\mathbb{F}}
\newcommand{\A}{\mathbb{A}}
\newcommand{\PP}{\mathbb{P}}
\newcommand{\E}{\mathbb{E}}
\newcommand{\eps}{\varepsilon}
\newcommand{\dd}{\,\mathrm{d}}
\begin{document}
\section*{UnsolvedMath ID: 30004608}
\textbf{Source problem code:} catalogue entry
\subsection{Definitions and mathematical statement}
The supplied question does not specify a ground field or precise conventions for plus-one generation and unexpected curves. The following gives an explicit conventional formulation over an algebraically closed field \(k\) of characteristic zero; this field hypothesis and the conventions below require confirmation against the source.

Let \(S=k[x,y,z]\), and let \(\mathcal A\) be a finite arrangement of distinct projective lines in \(\PP^2_k\), with reduced defining polynomial \(f=\prod_{L\in\mathcal A}\ell_L\) of degree \(n=|\mathcal A|\). Its intersection lattice \(L(\mathcal A)\) is the poset of intersections of the corresponding central hyperplanes in \(k^3\), ordered by reverse inclusion. A condition is combinatorial if it depends only on the isomorphism class of this lattice. The dual point set is
\[
Z_{\mathcal A}=\{[a_L:b_L:c_L]\in(\PP^2_k)^\vee : L=V(a_Lx+b_Ly+c_L)\in\mathcal A\}.
\]

Write
\[
\operatorname{AR}(f)=\{(u,v,w)\in S^3:u f_x+v f_y+w f_z=0\},
\]
with grading given by the degrees of the coefficients. A conventional graded formulation of plus-one generation is that \(\operatorname{AR}(f)\) has a minimal graded free resolution
\[
0\longrightarrow S(-c-1)\longrightarrow
S(-a)\oplus S(-b)\oplus S(-c)
\longrightarrow\operatorname{AR}(f)\longrightarrow0,
\qquad a\le b\le c,\quad a+b=n.
\]
Here \(S(-r)_q=S_{q-r}\). In this convention the nearly free case is the special case \(b=c\). These definitions supply a framework for the question; the record supplies neither the particular nearly free criterion to be generalized nor its hypotheses.

For a finite reduced point set \(Z\) in a projective plane, let \(\mathcal I_Z\) denote its ideal sheaf. Under the usual degree--multiplicity convention, \(Z\) admits an unexpected curve if there is an integer \(d\ge2\) such that, for a general point \(P\) of that plane,
\[
 h^0\!\left(\PP^2_k,
 (\mathcal I_Z\cap\mathcal I_P^{\,d-1})(d)\right)
 >\max\!\left\{0,
 h^0\!\left(\PP^2_k,\mathcal I_Z(d)\right)
 -\binom{d}{2}\right\}.
\]
Thus the curves have degree \(d\), contain \(Z\), and have multiplicity at least \(d-1\) at \(P\); their vector space has dimension greater than the expected dimension. Here general means that the inequality holds on a nonempty Zariski open set of choices of \(P\). Apply this definition in the dual plane for \(Z=Z_{\mathcal A}\). The source does not explicitly state whether this degree--multiplicity convention is intended.

\textbf{Problem.} Find explicit combinatorial sufficient conditions \(\mathcal C\), analogous to the conditions for nearly free arrangements intended by the source, such that every plus-one generated line arrangement \(\mathcal A\) satisfying those conditions has a dual point set admitting unexpected curves. In the provisional framework above, the desired implication is
\[
\forall\mathcal A,\qquad
\bigl(\mathcal A\text{ is plus-one generated}
\ \land\ \mathcal C(L(\mathcal A))\bigr)
\ \Longrightarrow\
\bigl(Z_{\mathcal A}\text{ admits an unexpected curve}\bigr).
\]
The source asks for sufficient criteria, without specifying a particular predicate \(\mathcal C\), a curve degree, or the nearly free result serving as the model. Consequently this is a research request for criteria, rather than a fully specified conjectural implication.
\subsection{Short English statement}
Find conditions determined by the intersection lattice of a plus-one generated line arrangement that guarantee unexpected curves through its dual points, in analogy with criteria for nearly free arrangements. The ground field, the precise conventions, and the nearly free criteria to be generalized are not specified in the supplied record. The displayed formulation uses explicit provisional conventions and requires human mathematical review.
\subsection{Sources}
\begin{itemize}
\item Local UnsolvedMath registry, supplied source record: permanent ID \(30004608\), title \emph{Unexpected Curves from Plus-One Generated Arrangements}, exact source problem code \texttt{catalogue entry}, collection \emph{Oberwolfach Reports Open Problems}. The record gives proposed year \(2021\) and names no individual proposer.
\item Supplied source citation: \emph{Logarithmic Vector Fields and Freeness of Divisors and Arrangements: New perspectives and applications} (2022), Oberwolfach Reports, pp.\ 225--300, DOI \texttt{10.4171/owr/2021/5}. \url{https://doi.org/10.4171/owr/2021/5}. The supplied background describes Oberwolfach Reports as documenting mathematical workshops and their open-problem discussions.
\item Supplied literature source: Bromboszcz, \emph{Rendiconti del Seminario Matematico della Universita di Padova}. \url{https://doi.org/10.1007/s40574-025-00517-z}. No article title or fuller bibliographic details are supplied.
\item Supplied literature source: \emph{Oberwolfach Report 5/2021}. \url{https://ems.press/content/serial-article-files/46883}.
\item Literature assessment dated exactly \texttt{2026-08-22}: the supplied curation labels the problem open and reports that recent work derives combinatorial constraints for plus-one generated line arrangements but does not establish the proposed general implication to unexpected curves. This is supplied dated curation, not an independently verified literature assessment. The background explicitly requests independent verification of time-sensitive status and cited claims.
\item The supplied statement status is \texttt{corrected\_verified}; its accompanying note reports a manual semantic, typographic, and LaTeX review of the clean statement. This is a supplied review label, not independent verification of the formulation, definitions, or mathematical status in this draft. No cited external source was visited in preparing this statement.
\end{itemize}
\end{document}
